\begin{proposition}[2.3.6]
\label{IV.2.3.6-fr}
Sea $Y$ un preesquema cuyo conjunto de componentes irreducibles sea
localmente finito \textup{(lo que ocurre si el espacio subyacente a $Y$ es
\emph{localmente noetheriano} (cf. \textbf{I}, 6.1.9))}.
\begin{enumerate}
  \item[\textup{(i)}] Toda parte cerrada $W$ de $Y$ estable por generización
  $(\mathbf{0}_{\mathrm I},2.1.2)$ es abierta. En particular, toda componente conexa
  de $Y$ es abierta.
  \item[\textup{(ii)}] Sea $f:X\to Y$ un morfismo \emph{cerrado} que satisfaga además las
  condiciones equivalentes \textup{(i), (ii), (iii)} de \textup{(2.3.4) (lo que ocurrirá si
  $f$ es \emph{cuasiplano} o \emph{abierto})}. Entonces la imagen por $f$ de toda componente
  conexa $C$ de $X$ es una componente conexa de $Y$.
  \item[\textup{(iii)}] Sea $f:X\to Y$ un morfismo que satisfaga las condiciones
  equivalentes \textup{(i), (ii), (iii)} de \textup{(2.3.4)} y tal, además, que para todo
  $y\in Y$, el conjunto $f^{-1}(y)$ sea finito \textup{(lo que ocurrirá si $f$ es
  \emph{cuasifinito} o \emph{radicial})}. Entonces el conjunto de componentes irreducibles de
  $X$ es localmente finito.
\end{enumerate}
\end{proposition}

(i) Si $y\in W$, los puntos genéricos $\eta_i$ ($1\leq i\leq m$) de las
componentes irreducibles $Y_i$ de $Y$ que contienen a $y$ pertenecen, por
hipótesis, a $W$; por tanto, puesto que $W$ es cerrado, esas mismas componentes
están contenidas en $W$; como, por hipótesis, existe una vecindad $U$ de $y$ tal
que $U$ es unión de los $U\cap Y_i$ $(\mathbf{0}_{\mathrm I},2.1.6)$, se tiene
$U\subset W$, por lo que $W$
\oldpage[IV-2]{16}
es abierto. Como, para todo $y\in Y$,
$\overline{\{y\}}$ es conexo, una componente conexa de $Y$ es estable por
generización; la segunda afirmación se deduce, pues, inmediatamente de la
primera.

(ii) Como $C$ es cerrado en $X$, $f(C)$ es cerrado en $Y$ por hipótesis.
Además, como $C$ es estable por generización, la hipótesis sobre $f$ implica
que $f(C)$ es estable por generización; por tanto, en virtud de (i), $f(C)$ es
abierto y cerrado y, como es conexo, es una componente conexa de $Y$.

(iii) Sea $x\in X$; por hipótesis, existe una vecindad abierta $U$ de
$y=f(x)$ que solo corta un número finito de componentes irreducibles $Y_i$ de
$Y$ ($1\leq i\leq n$); sea $y_i$ el punto genérico de $Y_i$
($1\leq i\leq n$). Para toda componente irreducible $Z$ de $X$ que corte a
$f^{-1}(U)$, el punto genérico $z$ de $Z$ pertenece necesariamente a uno de
los conjuntos $f^{-1}(y_i)$ (2.3.4). Como cada uno de esos conjuntos es finito
por hipótesis, queda demostrada nuestra afirmación.

\begin{proposition}[2.3.7]
\label{IV.2.3.7-fr}
Sean $f:X\to Y$ un morfismo, $g:Y'\to Y$ un morfismo cuasiplano,
$X'=X\times_Y Y'$, $f'=f_{(Y')}:X'\to Y'$.
\begin{enumerate}
  \item[\textup{(i)}] Si $f$ es cuasicompacto y dominante, también lo es $f'$.
  \item[\textup{(ii)}] Si toda componente irreducible de $X$ domina una componente
  irreducible de $Y$, entonces toda componente irreducible de $X'$ domina una
  componente irreducible de $Y'$.
\end{enumerate}
\end{proposition}

Designemos por $g':X'\to X$ la proyección canónica, que es un morfismo
cuasiplano (2.3.3).

(i) Ya se sabe (1.1.2) que $f'$ es cuasicompacto; además, si $y'$ es un punto
maximal (1.1.4) de $Y'$, $y=g(y')$ es un punto maximal de $Y$, como se deduce
de (2.3.4, (iii)). Por hipótesis (1.1.5), existe $x\in X$ tal que
$f(x)=y$; por tanto (\textbf I, 3.4.7), existe $x'\in X'$ tal que $f'(x')=y'$.

(ii) Sea $x'$ un punto maximal de $X'$, y sea $x=g'(x')$; se deduce de
(2.3.4, (ii)) que $x$ es un punto maximal de $X$, y por hipótesis $y=f(x)$ es
un punto maximal de $Y$. Pongamos $y'=f'(x')$, de modo que $g(y')=y$; se trata
de demostrar que $y'$ es un punto maximal de $Y'$. En virtud de (\textbf I, 5.1.7) y
(2.3.3), es posible limitarse al caso en que $X$ e $Y$ son \emph{reducidos}, por
lo que $\mathcal O_x$ y $\mathcal O_y$ son \emph{cuerpos}; además, en virtud
de (\textbf I, 3.6.5), aplicado dos veces, y de (\textbf I, 2.4.4), es posible limitarse al caso
en que $Y=\operatorname{Spec}(\mathcal O_y)$ y
$Y'=\operatorname{Spec}(\mathcal O_{y'})$. Entonces $f$ es un morfismo
\emph{plano}, puesto que $\mathcal O_y$ es un cuerpo (2.1.2), y lo mismo
ocurre con $f'$ (2.1.4); se deduce, pues, de (2.3.4, (ii)) que $y'=f'(x')$ es
un punto maximal de $Y'$.
